% Part: model-theory
% Chapter: basics
% Section: nonstandard-arithmetic

\documentclass[../../../include/open-logic-section]{subfiles}

\begin{document}

\olfileid{mod}{bas}{nsa}
\olsection{نماذج الحساب غير القياسية}

\begin{defn}
لتكن~$\Lang{L_N}$ !!{language} الحساب، وهي تتألف من~!!{constant}
$\Obj{0}$، و!!{predicate} ثنائي الموضع~$<$، و!!{function} أحادي الموضع
$\prime$، و!!{function}s ثنائية الموضع~$+$ و$\times$.
\begin{enumerate}
\item \emph{النموذج القياسي} للحساب هو~!!{structure}~$\Struct{N}$ للغة
  $\Lang{L_N}$ الذي مجاله~$\Nat = \{ 0, 1, 2, \dots\}$، ويفسر
  $\Obj{0}$ بالصفر~$0$، و$<$ بعلاقة الأصغر من على~$\Nat$، ويفسر
  $\prime$ و$+$ و$\times$ بالخلف والجمع والضرب على~$\Nat$ على
  الترتيب.
\item \emph{الحساب الصادق} هو النظرية~$\Theory{N}$.
\end{enumerate}
\end{defn}

حين نعمل في~$\Lang{L_N}$، نختصر كل حد من الصورة
$\Obj{0}^{\prime\dots\prime}$، وفيه~$n$ تطبيقات لدالة الخلف على
$\Obj{0}$، إلى~$\num{n}$.

\begin{defn}
يكون~!!{structure}~$\Struct{M}$ للغة~$\Lang{L_N}$ \emph{قياسيًا} إذا
وفقط إذا كان~$\Struct{N} \iso \Struct{M}$.
\end{defn}

\begin{thm}\ollabel{thm:non-std}
توجد نماذج~!!{enumerable} غير قياسية للحساب الصادق.
\end{thm}

\begin{proof}
وسّع~$\Lang{L_N}$ بإدخال~!!{constant} جديد~$c$، وتأمل النظرية
\[
\Theory{N} \cup \Setabs{\num{n} < c}{n \in \Nat}.
\]
هذه النظرية قابلة للإشباع منتهيًا، ولذلك لها بالتراص نموذج~$\Struct{M}$
يمكن اختياره~!!{enumerable} بمبرهنة لوفنهايم--سكولم النزولية. إذا كان
$\Domain{M}$ مجال~$\Struct{M}$، فليفسر~$\Struct{M}$ الرموز غير
المنطقية للغة~$\Lang{L_N}$ على النحو الآتي:
$\mathbf{z} = \Assign{\Obj{0}}{M} \in \Domain{M}$، و
${\prec} = \Assign{<}{M} \subseteq \Domain{M}^2$، و
$* = \Assign{\prime}{M} \colon \Domain{M} \to \Domain{M}$، و
$\oplus = \Assign{+}{M}$، و
$\otimes = \Assign{\times}{M} \colon \Domain{M}^2 \to \Domain{M}$.
ولكل~$x \in \Domain{M}$، نكتب~$x^*$ للعنصر من~$\Domain{M}$ الناتج من
تطبيق~$*$ على~$x$.

والآن، لو كان~$h$ تماثلًا بين~$\Struct{N}$ و$\Struct{M}$، لوجد
$n \in \Nat$ بحيث~$h(n) = \Assign{c}{M}$. فليكن~$s$ أي إسناد في
$\Struct{N}$ بحيث~$s(x) = n$. عندئذ
$\Sat{N}{\eq[\num{n}][x]}[s]$؛ وببرهان~\olref[iso]{thm:isom} يكون
أيضًا~$\Sat{M}{\eq[\num{n}][x]}[h \circ s]$، ومن ثم
$\Assign{c}{M} = \mathbf{z}^{*\cdots *}$، حيث تكرر~$*$ عدد~$n$ من
المرات. لكن ذلك مستحيل، لأن الفرض يعطي
$\Sat{M}{\num{n} < c}$، ولأن~$\prec$ لا انعكاسية. ولذلك يكون
$\Struct{M}$ غير قياسي.
\end{proof}

\begin{prob}
تكون علاقة~$R$ على مجموعة~$X$ \emph{راسخة} إذا وفقط إذا لم توجد فيها
سلاسل نزولية غير منتهية؛ أي إذا لم توجد~$x_0$،~$x_1$،~$x_2$،~\dots
في~$X$ بحيث~$\dots x_2Rx_1Rx_0$. وبافتراض اتساق نظرية مجموعات
زيرميلو--فرينكل~$ZF$، بيّن أن لـ$ZF$ نماذج غير راسخة؛ أي نماذج
$\mathfrak{M}$ بحيث~$\dots x_2 \in x_1 \in x_0$.
\end{prob}

لما كان النموذج غير القياسي~$\Struct{M}$ من~\olref{thm:non-std}
مكافئًا أوليًا للنموذج القياسي، أمكن اشتقاق عدد من خواص~$\Struct{M}$.
ويُكرس باقي هذا القسم لهذه المهمة، مما يتيح لنا وصف نماذج~$\Theory{N}$
غير القياسية~!!{enumerable} وصفًا دقيقًا.

\begin{enumerate}
\item لا يكون أي عنصر من~$\Domain{M}$ أصغر بحسب~$\prec$ من نفسه:
  فالجملة~$\lforall[x][\lnot x < x]$ صادقة في~$\Struct{N}$، ومن ثم
  في~$\Struct{M}$.
\item وباستدلال مماثل نحصل على أن~$\prec$ \emph{ترتيب خطي} لـ
  $\Domain{M}$، أي علاقة كلية لا انعكاسية ومتعدية على~$\Domain{M}$.
\item العنصر~$\mathbf{z}$ هو أصغر عناصر~$\Domain{M}$ بحسب~$\prec$.
\item كل عنصر من~$\Domain{M}$ أصغر بحسب~$\prec$ من خلفه بالنسبة إلى
  $*$، و$x^*$ هو أصغر عناصر~$\Domain{M}$ بحسب~$\prec$ بين العناصر
  الأكبر من~$x$.
\item يحتوي~$\Struct{M}$ قطعة ابتدائية بالنسبة إلى~$\prec$ متماثلة
  مع~$\Nat$، هي~$\mathbf{z}, \mathbf{z}^*, \mathbf{z}^{**}, \dots$،
  ونسميها~\emph{الجزء القياسي} من~$\Domain{M}$. وكل عنصر آخر من
  $\Domain{M}$ \emph{غير قياسي}. ولا بد من وجود عناصر غير قياسية في
  $\Domain{M}$، وإلا لكانت الدالة~$h$ من برهان~\olref{thm:non-std}
  تماثلًا. ونستخدم~$n, m, \dots$ بوصفها~!!{variable}s تتراوح على هذا
  الجزء القياسي من~$\Struct{M}$.
\item كل عنصر غير قياسي أكبر من كل عنصر قياسي؛ ذلك أنه لكل
  $n \in \Nat$،
  \[
  \Sat{N}{\lforall[z][(\lnot(\eq[z][\Obj{0}] \lor
  \dots \lor \eq[z][\num{n}]) \lif \num{n} < z)]},
  \]
  ولذلك إذا كان~$z \in \Domain{M}$ مختلفًا عن جميع العناصر القياسية،
  وجب أن يكون~\emph{أكبر} منها جميعًا.
\item كل عنصر من~$\Domain{M}$ غير~$\mathbf{z}$ هو خلف بالنسبة إلى~$*$
  لعنصر وحيد من~$\Domain{M}$، نرمز إليه بـ$^*x$. وإذا كان~$x = y^*$،
  كان~$x$ و$y$ قياسيين إذا كان أحدهما قياسيًا، وكانا غير قياسيين إذا
  كان أحدهما غير قياسي.
\item عرّف علاقة تكافؤ~$\approx$ على~$\Domain{M}$ بالقول إن
  $x \approx y$ إذا وفقط إذا كان، لبعض~$n$ \emph{قياسي}، إما
  $x \oplus n = y$ وإما~$y \oplus n =x$. وبعبارة أخرى، يكون
  $x \approx y$ إذا وفقط إذا كانت المسافة بين~$x$ و$y$ منتهية. وإذا
  كان~$n$ و$m$ قياسيين، كان~$n \approx m$. وعرّف~\emph{كتلة}~$x$
  بأنها صنف التكافؤ~$[x] = \Setabs{y}{x \approx y}$.
\item افترض~$x \prec y$ و$x \not\approx y$. ولأن
  $\Sat{N}{\lforall[x][\lforall[y][(x < y \lif (x' < y \lor x' =
      y))]]}$، يكون إما~$x^* \prec y$ وإما~$x^* = y$. والحالة
  الأخيرة مستحيلة لأنها تستلزم~$x \approx y$، ومن ثم~$x^* \prec y$.
  وبالمثل، إذا كان~$x \prec y$ و$x \not\approx y$، كان
  $x \prec {^*y}$. لذلك، إذا كان~$x \prec y$ و$x \not\approx y$،
  كان كل~$w \approx x$ أصغر بحسب~$\prec$ من كل~$v \approx y$.
  وكل كتلة غير قياسية~$[x]$ تكوّن سلسلة غير منتهية في الاتجاهين
  \[
  \cdots \prec {^{**}x} \prec {^*}x \prec x \prec x^* \prec x^{**}
  \prec \cdots
  \]
  تسمى سلسلة~$Z$ لأن نوع ترتيبها هو نوع ترتيب الأعداد الصحيحة.
\item يمكن رفع ترتيب~$\prec$ إلى الكتل: إذا كان~$x \prec y$ و
  $x \not\approx y$، كانت كتلة~$x$ أصغر من كتلة~$y$. وتكون الكتلة
  \emph{غير قياسية} إذا احتوت عنصرًا غير قياسي. أما الكتلة القياسية
  فهي أصغر الكتل.
\item لا توجد أصغر كتلة غير قياسية: فإذا كان~$y$ غير قياسي، وُجد
  $x \prec y$ يكون~$x$ فيه غير قياسي أيضًا و$x \not\approx y$.
  البرهان: في النموذج القياسي~$\Struct{N}$، يقبل كل عدد القسمة على
  اثنين، وربما كان الباقي واحدًا؛ أي
  $\Sat{N}{\lforall[y][\lexists[x][(y = x + x \lor y = x + x +
        \Obj{0}')]]}$. وبالتكافؤ الأولي، لكل~$y \in \Domain{M}$ يوجد
  $x \in \Domain{M}$ بحيث إما~$x \oplus x = y$ وإما
  $x \oplus x \oplus \mathbf{z}^*= y$. ولو كان~$x$ قياسيًا، لكان~$y$
  قياسيًا أيضًا؛ ولذلك~$x$ غير قياسي، كما أن~$x \prec y$. وفوق ذلك،
  ينتمي~$x$ و$y$ إلى كتلتين مختلفتين، أي~$x \not\approx y$. ولرؤية
  ذلك، افترض أنهما ينتميان إلى الكتلة نفسها، أي
  $x \oplus n = y$ لبعض~$n$ قياسي. فإذا كان~$y = x \oplus x$، كان
  $x \oplus n = x \oplus x$، ومنه~$x = n$ بقانون الاختزال للجمع،
  وهو قانون يصدق في~$\Struct{N}$ ومن ثم في~$\Struct{M}$ أيضًا، فيكون
  $x$ قياسيًا في نهاية المطاف. وتُعالج حالة
  $y = x \oplus x \oplus \mathbf{z}^*$ على نحو مماثل.
\item وبحجة مماثلة، لا توجد أكبر كتلة.
\item ترتيب الكتل كثيف: إذا كانت~$[x]$ أصغر من~$[y]$، حيث
  $x \not\approx y$، وُجدت كتلة~$[z]$ متميزة منهما وتقع بينهما. افترض
  $x \prec y$. وكما سبق، يقبل~$x \oplus y$ القسمة على اثنين، وربما
  كان ثمة باق، ولذلك يوجد~$u \in \Domain{M}$ بحيث إما
  $x \oplus y = u \oplus u$ وإما
  $x \oplus y = u \oplus u \oplus \mathbf{z}^*$. والعنصر~$u$ هو
  متوسط~$x$ و$y$، ولذلك يقع بينهما. افترض
  $x \oplus y = u \oplus u$، والحالة الأخرى مماثلة: إذا كان
  $u \approx x$، فلبعض~$n$ قياسي
  \[
  x \oplus y = x \oplus n \oplus x \oplus n,
  \]
  ومن ثم~$y = x \oplus n \oplus n$، فيكون~$x \approx y$ خلافًا
  للفرض. فنستنتج~$u \not\approx x$. وتعطي حجة مماثلة
  $u \not\approx y$.
\end{enumerate}
وهكذا تُرتب الكتل غير القياسية مثل الأعداد النسبية: فهي تكوّن ترتيبًا
خطيًا~!!a{enumerable} بلا نقطتي نهاية. ويلزم من ذلك أنه لأي نموذجين
غير قياسيين~!!{enumerable}~$\Struct{M}_1$ و$\Struct{M}_2$ للحساب
الصادق، يكون اختزالهما إلى اللغة التي لا تحتوي إلا~$<$ و$=$ متماثلين.
فالجزآن القياسيان من~$\Struct{M_1}$ و$\Struct{M_2}$ متماثلان مع
النموذج القياسي~$\Struct{N}$، ومن ثم أحدهما مع الآخر. والكتل المكوّنة
للجزء غير القياسي مرتبة هي نفسها مثل الأعداد النسبية، ولذلك تكون
متماثلة بموجب~\olref[dlo]{thm:cantorQ}؛ ويمكن توسيع تماثل الكتل إلى
تماثل~\emph{داخل} الكتل بمقابلة عناصر اعتباطية في كل منها، ثم جعل
صورة خلف~$x$ في~$\Struct{M_1}$ خلف صورة~$x$ في~$\Struct{M_2}$. لاحظ
أنه لا يلزم من ذلك~\emph{أن}~$\mathfrak{M}_1$ و$\mathfrak{M}_2$
متماثلان في لغة الحساب الكاملة، فالتماثل دائمًا نسبي إلى توقيع؛ إذ
توجد طرائق غير متماثلة لتعريف الجمع والضرب على~$\Domain{M_1}$ و
$\Domain{M_2}$. (ويلزم ذلك أيضًا من مبرهنة شهيرة لفوت، تقول إن عدد
النماذج المعدودة لنظرية تامة لا يمكن أن يكون~2.)

\begin{prob}
بيّن أنه لا يمكن أن توجد أكبر كتلة في نموذج حسابي غير قياسي.
\end{prob}

\begin{prob}
لتكن~$\Lang{L}$ !!{language} الرتبة الأولى التي لا تحتوي، عدا~$\eq$،
إلا~$<$ بوصفه~!!{predicate} وحيدًا، ولتكن
$\Struct{N} = (\Nat,<)$. جميع المجموعات الجزئية المنتهية أو المتناهية
المتممة من~$\Struct{N}$ قابلة للتعريف من دون معاملات. بيّن أنها
المجموعات الجزئية الوحيدة القابلة للتعريف من دون معاملات في
$\Struct{N}$.

(تلميح: أولًا، لتكن~$\Obj{prc}(x,y)$ صيغة~$\Lang{L}$ المختصرة لعبارة
«$x$ هو السلف المباشر لـ$y$»:
\[
x<y \land \lnot \lexists[z][(x<z \land z < y)].
\]
تقابل كل مجموعة جزئية من~$\Struct{N}$ قابلة للتعريف من دون معاملات
صيغة~$!A(x)$ في~$\Lang{L}$. ولكل صيغة كهذه~$!A$، تأمل الجملة~$!D$:
\[
\lexists[x][\lforall[y][\lforall[z][((x<y \land x<z \land \Obj{prc}(y,z)
  \land !A(y)) \lif !A(z))]]].
\]
بيّن أن~$\Sat{N}{!D}$ إذا وفقط إذا كانت المجموعة الجزئية من
$\Struct{N}$ التي تعرّفها~$!A$ إما منتهية وإما متناهية المتممة.

والآن ليكن~$\Struct{M}$ نموذجًا غير قياسي مكافئًا أوليًا لـ
$\Struct{N}$. وإذا كان~$a \in \Domain{M}$ غير قياسي، فليكن
$b, c \in \Domain{M}$ أكبر من~$a$، وليكن~$b$ السلف المباشر لـ$c$.
عندئذ يوجد تشاكل ذاتي~$h$ لاختزال~$\Struct{M}$ إلى لغة الترتيب بحيث
$h(b)=c$، فلماذا؟ ولذلك إذا أشبع~$b$ الصيغة~$!A$، أشبعها~$c$ أيضًا،
فلماذا؟ ويلزم أن~$!D$ صادقة في~$\Struct{M}$، ومن ثم في~$\Struct{N}$
أيضًا. لكن ذلك يستلزم أن المجموعة الجزئية من~$\Struct{N}$ التي
تعرّفها~$!A$ إما منتهية وإما متناهية المتممة.)
\end{prob}

\end{document}
